\begin{lemma}[The $\varepsilon$-independent majorant for the truncated Mehler gradient slice]
Fix $n \geq 1$, $j \in \{1,\ldots,n\}$ and $x \in \mathbb{R}^n$.  For $s > 0$
write
\[
  r = e^{-s} \in (0,1), \qquad \sigma = 1 - e^{-2s} \in (0,1),
  \qquad
  p_s(x,y) = \frac{1}{\pi^{n/2}\sigma^{n/2}}
      \exp\Bigl(-\frac{|rx-y|^2}{\sigma}\Bigr) ,
\]
and let
\[
  D_j^x p_s(x,y) = \sqrt2\,r\,\frac{y_j - r x_j}{\sigma}\,p_s(x,y)
\]
be the energy-normalized $x$-derivative of the Lebesgue kernel, as evaluated in
\noderef{MehlerKernelXDerivative} (the displayed right-hand side is the
closed-form expression that node proves $2^{-1/2}\partial_{x_j}p_s(x,y)$ equals,
and it is that explicit expression, not a derivative, that occurs here).  For
$\varepsilon > 0$ put
$\Omega_\varepsilon = \{\,y \in \mathbb{R}^n : |x-y| > \varepsilon\,\}$ and
\[
  I_{\varepsilon,s} = \int_{\Omega_\varepsilon} D_j^xp_s(x,y)\,dy .
\]
Then:
\begin{enumerate}
\item there is a constant $C \geq 0$ such that
\[
  \bigl|I_{\varepsilon,s}\bigr| \leq C\,|x|\,e^{-s}
  \qquad\text{for every } \varepsilon > 0 \text{ and every } s > 0 ;
\]
\item for every $s > 0$,
$\displaystyle\lim_{\varepsilon\downarrow0} I_{\varepsilon,s} = 0$, the limit
being taken along $\varepsilon \to 0$ with $\varepsilon > 0$.
\end{enumerate}

The force of clause (1) is the quantifier order: $C$ is chosen after $n$, $j$
and $x$ and so may depend on all three, but \emph{before} $\varepsilon$ and $s$,
so it is independent of both.  That independence is the whole point --- the
truncated integrals blow up like $\sigma^{-1/2}$ if the cancellation between the
two translates is not used, and the bound above is uniform in the truncation
parameter and decays in $s$.  This is the constant the reference paper's
appendix header licenses: it may depend on $n$, it serves only to identify two
realizations of one operator, and no claim of uniformity in $n$ is made or used
here.  Clause (2) is the fixed-time limit, with no claim of uniformity in $s$.
Nothing is asserted about the normalized Gaussian Riesz kernel, about any time
integral, or about any exchange of the two integrations.
\end{lemma}
\begin{proof}
Throughout, $\lambda$ denotes Lebesgue measure on $\mathbb{R}^n$, and $n \geq 1$,
$j$ and $x$ are fixed.  We take
\[
  C = \sqrt2\,\bigl(5 + 24\cdot8^{n/2}\bigr) \geq 0 ,
\]
which depends on $n$ only.

\emph{Step 0: a Gaussian toolkit.}
First, for every real $u \geq 0$,
\begin{equation*}
  u\,e^{-u} \leq 1 . \tag{0}
\end{equation*}
Indeed $1+u \leq e^u$ (shared context, the exponential inequality), so
$u \leq e^u$, and multiplying by $e^{-u} > 0$ gives $(0)$.

Second, for every $\beta > 0$ the function $w \mapsto e^{-\beta|w|^2}$ is
Lebesgue integrable on $\mathbb{R}^n$ with
\begin{equation*}
  \int_{\mathbb{R}^n} e^{-\beta|w|^2}\,dw = (\pi/\beta)^{n/2} . \tag{1}
\end{equation*}
Indeed $e^{-\beta|w|^2} = \prod_{i=1}^n e^{-\beta w_i^2}$, Lebesgue measure on
$\mathbb{R}^n$ is the $n$-fold product of one-dimensional Lebesgue measure, and
the integrand is nonnegative and continuous, hence product-measurable; so
Tonelli (shared context) and the scalar Gaussian integral (shared context,
$\int_{\mathbb{R}}e^{-\beta v^2}dv = \sqrt{\pi/\beta}$) give $(1)$, the value
being finite.

Third, for every $\beta > 0$ and every $w$,
\begin{equation*}
  |w|^2e^{-\beta|w|^2} \leq \frac2\beta\,e^{-(\beta/2)|w|^2} , \tag{2}
\end{equation*}
whence $w \mapsto |w|^2e^{-\beta|w|^2}$ is integrable, being continuous and
dominated by an integrable function (shared context), and
\begin{equation*}
  \int_{\mathbb{R}^n}|w|^2e^{-\beta|w|^2}\,dw
    \leq \frac2\beta\,(2\pi/\beta)^{n/2} . \tag{3}
\end{equation*}
For $(2)$: put $u = (\beta/2)|w|^2 \geq 0$; then
$e^{-\beta|w|^2} = e^{-u}e^{-(\beta/2)|w|^2}$, and $(0)$ gives
$u\,e^{-u} \leq 1$, i.e.\ $|w|^2e^{-u} \leq 2/\beta$ after multiplying by
$2/\beta > 0$; multiplying by $e^{-(\beta/2)|w|^2} > 0$ gives $(2)$.  Then $(3)$
follows by integrating $(2)$ and using $(1)$ at $\beta/2$.

Fourth, for all reals $a$, $b$, $m$ with $m \leq a$ and $m \leq b$,
\begin{equation*}
  \bigl|e^{-a}-e^{-b}\bigr| \leq |a-b|\,e^{-m} . \tag{4}
\end{equation*}
By symmetry assume $a \leq b$.  Then $e^{-b} \leq e^{-a}$, so the left side is
$e^{-a}-e^{-b} = e^{-a}\bigl(1-e^{-(b-a)}\bigr)$; and $1-e^{-t} \leq t$ for
$t \geq 0$ by the exponential inequality again, with $t = b-a \geq 0$; so the
left side is at most $e^{-a}(b-a) \leq e^{-m}|a-b|$, using $e^{-a} \leq e^{-m}$.

Fifth, $|w_j| \leq |w|$ for every $w$: writing
$|w| = (\sum_i|w_i|^2)^{1/2}$, a sum of nonnegative terms is at least its $j$th
term, so $|w_j|^2 \leq |w|^2$, and the square root is monotone.

\emph{Step 1: the integrand after translation.}
Fix $s > 0$ and put $\beta = 1/\sigma > 0$,
\[
  c = \frac{\sqrt2\,r\,\pi^{-n/2}\sigma^{-n/2}}{\sigma} \geq 0,
  \qquad d = (1-r)x,
  \qquad W(w) = c\,w_j\,e^{-\beta|w|^2} .
\]
For every $w$ we have $(rx+w)_j - rx_j = w_j$ and $rx - (rx+w) = -w$, so
$|rx-(rx+w)| = |w|$ and hence
\begin{equation*}
  D_j^xp_s(x,\,rx+w) = W(w) . \tag{5}
\end{equation*}
Note $|d| = (1-r)|x|$ since $1-r > 0$, and $\sigma = 1-r^2 = (1-r)(1+r)$.

\emph{Step 2: the pointwise majorant for the translation difference.}
We claim that for every $z$,
\begin{equation*}
  \bigl|W(z+d)-W(z)\bigr|
    \leq c\,|d|\Bigl(3e^{-\beta|z+d|^2} + 2e^{-\beta|z|^2}
        + 3\beta\,|z|^2e^{-(\beta/4)|z|^2}\Bigr) =: M(z) . \tag{6}
\end{equation*}
Write $E(w) = e^{-\beta|w|^2}$.  Since $(z+d)_j = z_j + d_j$,
\[
  W(z+d)-W(z) = c\,d_j\,E(z+d) + c\,z_j\bigl(E(z+d)-E(z)\bigr) ,
\]
so by the triangle inequality, $|d_j| \leq |d|$ and $E > 0$,
\begin{equation*}
  \bigl|W(z+d)-W(z)\bigr|
    \leq c\,|d|\,E(z+d) + c\,|z_j|\,\bigl|E(z+d)-E(z)\bigr| . \tag{7}
\end{equation*}
If $|z| \leq 2|d|$, then $|z_j| \leq |z| \leq 2|d|$ and
$|E(z+d)-E(z)| \leq E(z+d)+E(z)$, so the second term of $(7)$ is at most
$2c|d|(E(z+d)+E(z))$ and the right side of $(7)$ is at most
$c|d|(3E(z+d)+2E(z))$, which is at most $M(z)$ because the remaining term of
$M$ is nonnegative.

If $|z| > 2|d|$, put $a = \beta|z+d|^2$, $b = \beta|z|^2$ and
$m = (\beta/4)|z|^2$.  Then $m \leq b$; and $|d| \leq |z|/2$ gives
$|z+d| \geq |z|-|d| \geq |z|/2$, so $|z+d|^2 \geq |z|^2/4$ and $m \leq a$.
Moreover $|z+d|^2 - |z|^2 = 2\langle z,d\rangle + |d|^2$, so by Cauchy--Schwarz
(shared context) and $|d|^2 \leq (|z|/2)|d|$,
\[
  |a-b| = \beta\bigl|2\langle z,d\rangle+|d|^2\bigr|
    \leq \beta\bigl(2|z||d| + \tfrac12|z||d|\bigr) \leq 3\beta|z||d| .
\]
By $(4)$, $|E(z+d)-E(z)| \leq 3\beta|z||d|\,e^{-(\beta/4)|z|^2}$, so with
$|z_j| \leq |z|$ the second term of $(7)$ is at most
$3\beta c|d|\,|z|^2e^{-(\beta/4)|z|^2}$, and the right side of $(7)$ is at most
$c|d|\bigl(E(z+d) + 3\beta|z|^2e^{-(\beta/4)|z|^2}\bigr) \leq M(z)$.  This
proves $(6)$ in both cases.

\emph{Step 3: integrating the majorant.}
By $(1)$, translation invariance of $\lambda$ (shared context) and $(3)$ at
$\beta/4$, each of the three summands of $M$ is integrable and
\begin{align*}
  \int_{\mathbb{R}^n} M
    &= c|d|\Bigl(3(\pi/\beta)^{n/2} + 2(\pi/\beta)^{n/2}
        + 3\beta\int|z|^2e^{-(\beta/4)|z|^2}dz\Bigr)\\
    &\leq c|d|\Bigl(5(\pi/\beta)^{n/2}
        + 3\beta\cdot\frac8\beta\,(8\pi/\beta)^{n/2}\Bigr)
     = c|d|\,(\pi/\beta)^{n/2}\bigl(5+24\cdot8^{n/2}\bigr) ,
\end{align*}
using $(8\pi/\beta)^{n/2} = 8^{n/2}(\pi/\beta)^{n/2}$.  The function
$z \mapsto W(z+d)-W(z)$ is integrable: $W$ is integrable because, by $(5)$ and
translation invariance, it is the composition of
$y \mapsto D_j^xp_s(x,y)$ with a translation, and that function is integrable by
clause (1) of \noderef{MehlerKernelSliceVanishing}; and $z \mapsto W(z+d)$ is
then integrable by translation invariance again.  Hence, by $(6)$ and
monotonicity of the integral,
\begin{equation*}
  \int_{\mathbb{R}^n}\bigl|W(z+d)-W(z)\bigr|\,dz
    \leq c|d|\,(\pi/\beta)^{n/2}\bigl(5+24\cdot8^{n/2}\bigr) . \tag{8}
\end{equation*}
Steps 2 and 3 together are \noderef{GaussFirstMomentTranslationL1Bound}, which states
exactly the bound $(8)$ for every $b > 0$, every $c \geq 0$ and every $d$, and is what
the formalization invokes here; the Gaussian toolkit of Step 0 that it needs is in turn
supplied to it by \noderef{GaussEuclideanMomentBounds}.

\emph{Step 4: two translations and the oddness.}
Fix $\varepsilon > 0$.  The set $\Omega_\varepsilon$ is open, hence measurable,
$y \mapsto |x-y|$ being continuous.  Translating by $rx$: the preimage of
$\Omega_\varepsilon$ under $w \mapsto rx+w$ is
$\{w : |d-w| > \varepsilon\}$, because $x-(rx+w) = d-w$; so by translation
invariance of $\lambda$ and $(5)$,
\[
  I_{\varepsilon,s} = \int_{\{|d-w|>\varepsilon\}} W(w)\,dw .
\]
Translating by $d$: the preimage of $\{|d-w|>\varepsilon\}$ under
$z \mapsto z+d$ is $\{z : |z| > \varepsilon\}$, because $d-(z+d) = -z$ and
$|-z| = |z|$; so
\[
  I_{\varepsilon,s} = \int_{\{|z|>\varepsilon\}} W(z+d)\,dz .
\]
Next, $W$ is odd: $(-z)_j = -z_j$ and $|-z| = |z|$ give $W(-z) = -W(z)$.  The
set $\{|z|>\varepsilon\}$ is invariant under $z \mapsto -z$, again because
$|-z| = |z|$, and Lebesgue measure is invariant under that reflection (shared
context); so, $W$ being integrable,
\[
  \int_{\{|z|>\varepsilon\}} W
    = \int_{\{|z|>\varepsilon\}} W(-z)\,dz
    = -\int_{\{|z|>\varepsilon\}} W ,
\]
whence $\int_{\{|z|>\varepsilon\}}W = 0$.  Subtracting,
\[
  I_{\varepsilon,s}
    = \int_{\{|z|>\varepsilon\}}\bigl(W(z+d)-W(z)\bigr)\,dz ,
\]
both terms being integrable, and therefore, by $(8)$ and the fact that the
integral of a nonnegative integrable function over a measurable set is at most
its integral over $\mathbb{R}^n$,
\begin{equation*}
  \bigl|I_{\varepsilon,s}\bigr|
    \leq \int_{\mathbb{R}^n}\bigl|W(z+d)-W(z)\bigr|\,dz
    \leq c|d|\,(\pi/\beta)^{n/2}\bigl(5+24\cdot8^{n/2}\bigr) . \tag{9}
\end{equation*}

\emph{Step 5: the constant, and clause (1).}
Since $\beta = 1/\sigma$ we have $\pi/\beta = \pi\sigma$ and therefore
$(\pi/\beta)^{n/2} = \pi^{n/2}\sigma^{n/2}$, so
\[
  c\,(\pi/\beta)^{n/2}
    = \frac{\sqrt2\,r}{\sigma}\,
      \bigl(\pi^{-n/2}\pi^{n/2}\bigr)\bigl(\sigma^{-n/2}\sigma^{n/2}\bigr)
    = \frac{\sqrt2\,r}{\sigma} .
\]
Using $|d| = (1-r)|x|$ and $\sigma = (1-r)(1+r)$ with $1-r > 0$,
\[
  \frac{\sqrt2\,r}{\sigma}\,|d|
    = \sqrt2\,|x|\,r\cdot\frac{1}{1+r}
    \leq \sqrt2\,|x|\,r ,
\]
because $0 \leq 1/(1+r) \leq 1$ and $\sqrt2|x|r \geq 0$.  Substituting into
$(9)$,
\[
  \bigl|I_{\varepsilon,s}\bigr|
    \leq \sqrt2\,\bigl(5+24\cdot8^{n/2}\bigr)\,|x|\,e^{-s} = C\,|x|\,e^{-s} ,
\]
for every $\varepsilon > 0$ and every $s > 0$, which is clause (1).

\emph{Step 6: clause (2).}
Fix $s > 0$.  By clause (1) of \noderef{MehlerKernelSliceVanishing} the function
$y \mapsto D_j^xp_s(x,y)$ is integrable on $\mathbb{R}^n$ with integral $0$, so
splitting the integral over $\Omega_\varepsilon$ and its complement,
\[
  I_{\varepsilon,s}
    = -\int_{\Omega_\varepsilon^{\,c}} D_j^xp_s(x,y)\,dy
    = -\int_{\mathbb{R}^n}
        \mathbf 1_{\Omega_\varepsilon^{\,c}}(y)\,D_j^xp_s(x,y)\,dy ,
\]
where $\Omega_\varepsilon^{\,c} = \{y : |x-y| \leq \varepsilon\}$.  For each
$\varepsilon > 0$ the integrand is measurable, being the product of a
measurable indicator with an integrable function, and it is dominated in
absolute value by $|D_j^xp_s(x,\cdot)|$, which is integrable and independent of
$\varepsilon$.  Moreover, for every $y \neq x$ we have $|x-y| > 0$, and
$\{\varepsilon : \varepsilon < |x-y|\}$ is a neighbourhood of $0$, so for all
$\varepsilon$ in a neighbourhood of $0$ within $(0,\infty)$ we have
$\varepsilon < |x-y|$, i.e.\ $y \notin \Omega_\varepsilon^{\,c}$ and the
integrand vanishes at $y$; hence the integrand tends to $0$ as
$\varepsilon \downarrow 0$ for every such $y$.  Since $n \geq 1$, the singleton
$\{x\}$ is $\lambda$-null, so this holds $\lambda$-almost everywhere.  Dominated
convergence along the filter of neighbourhoods of $0$ within $(0,\infty)$
(shared context; that filter is countably generated, $\mathbb{R}$ being
first countable) therefore gives
\[
  \int_{\mathbb{R}^n}
      \mathbf 1_{\Omega_\varepsilon^{\,c}}(y)\,D_j^xp_s(x,y)\,dy
    \longrightarrow \int_{\mathbb{R}^n} 0\,dy = 0
  \qquad (\varepsilon\downarrow0) ,
\]
and hence $I_{\varepsilon,s} \to 0$, which is clause (2).
\end{proof}
